\documentclass[11pt]{article}

\usepackage[margin=1in]{geometry}
\usepackage{booktabs}
\usepackage{hyperref}
\usepackage{enumitem}
\usepackage{listings}
\usepackage{xcolor}

\hypersetup{
  colorlinks=true,
  linkcolor=blue,
  citecolor=blue,
  urlcolor=blue
}

\lstset{
  basicstyle=\ttfamily\small,
  columns=fullflexible,
  frame=single,
  breaklines=true,
  backgroundcolor=\color{black!3}
}

\title{\textbf{HATS: A Heterogeneous Agent Tool Substrate for Agentic Prefill}}
\author{HATS Project}
\date{\today}

\begin{document}
\maketitle

\begin{abstract}
Agentic AI systems increasingly combine large model inference with tool use:
retrieval, structured parsing, schema validation, ranking, memory lookup,
code-like transformations, cancellation, retry, and dependency scheduling.
Meanwhile, much accelerator work is specializing the autoregressive decode path:
one-token continuation, KV-cache reads, and even separate attention and feed
forward execution. HATS takes the complementary position. It specializes the
agentic prefill path: the construction, validation, retrieval, filtering, and
ingestion of context before model continuation. In today's systems, the dense
tensor work is placed on GPUs, while much of this agentic prefill and re-prefill
work remains on the CPU. This split is convenient, but it can create
high-latency, high-power execution patterns: model decoding emits a tool
request, host software wakes up, validates arguments, launches another kernel or
IO operation, copies or maps results, builds continuation context, and then
resumes GPU-side inference.

This paper proposes HATS, a Heterogeneous Agent Tool Substrate inspired by HSA
principles but scoped to bounded local tools for agentic AI. HATS is not a
proposal to replace the CPU with a GPU control processor. Instead, it is a
compiler/runtime substrate that exposes coherent memory, user-mode queues,
fine-grained signals, capability-scoped tool access, and a narrow set of
hardware-friendly tool primitives. The initial target set includes retrieval,
structured parsing, schema validation, context construction, tool-result
ingestion, dependency scheduling, data movement, and sandboxed pure compute
kernels.

The central claim is conditional and empirical: agentic AI does not justify
moving arbitrary external tools into accelerator hardware, but it may justify
moving frequent, memory-resident prefill-side tool primitives into a coherent
GPU-side runtime. HATS therefore begins as a compiler and runtime project that
can run on existing heterogeneous systems, measure bottlenecks, and identify
which hardware primitives have credible performance, power, and area value.
\end{abstract}

\section{Introduction}

Modern AI accelerators are designed around dense numerical kernels. This has
served neural network training and inference well: the central optimization
problem is to move tensors through high-throughput matrix engines efficiently.
Agentic AI changes the workload shape. A useful agent does not merely decode
tokens; it plans, branches, retrieves memory, validates structured arguments,
invokes tools, cancels stale work, merges results, ranks alternatives, builds
fresh context, and resumes model execution.

This paper focuses on the prefill side of that loop. Conventional prefill
converts an input prompt into model state, including a KV cache that can be used
by subsequent decode. Agentic systems create prefill-like events repeatedly:
initial prompts are prefetched once, but tool results, retrieved memories,
ranked candidates, and repaired structured arguments must be ingested throughout
the workflow. These re-prefill events are where HATS expects the strongest
opportunity. Decode should continue to be specialized aggressively, but HATS
targets the repeated construction of useful context around decode.

This outer loop is not a small detail. Many agentic workflows spend substantial
time outside the core matrix multiplication path. A common execution pattern is:

\begin{lstlisting}
GPU decode emits a tool boundary
  -> CPU observes tool request
  -> CPU validates arguments
  -> CPU launches retrieval or parsing work
  -> GPU or CPU computes the tool result
  -> CPU synchronizes and formats the result
  -> CPU/GPU builds continuation context and KV state
  -> GPU resumes decoding
\end{lstlisting}

Each transition may be reasonable in isolation. Together, the pattern can
produce poor locality, frequent host wakeups, heavyweight synchronization, and
unnecessary movement of data that was already close to GPU-resident model state.
The power and latency cost is especially visible when tool calls are small,
frequent, and dependent on fine-grained model outputs.

HATS asks whether the agentic outer loop should have a first-class
heterogeneous execution substrate. The answer should not be assumed. General
branching, OS policy, arbitrary network APIs, browsers, and shell tools remain
natural CPU responsibilities. The narrower opportunity is different: many
agentic systems repeatedly use bounded, local, memory-oriented primitives that
are close to the model execution path. These include vector retrieval over
resident indexes, key-value memory lookup, JSON or schema validation, constrained
decoding checks, fan-out/fan-in scheduling, cancellation propagation, and small
deterministic scoring or filtering functions.

HATS is inspired by HSA's core philosophy: heterogeneous processors should
cooperate through shared virtual memory, queues, and signals rather than through
heavy host-mediated boundaries. HATS applies that philosophy to agentic tool
workflows. The result is not a general-purpose operating system inside a GPU,
but a bounded substrate for local tool primitives in the agentic prefill path.

\section{Motivation}

\subsection{Agentic Workloads Are Not Pure Tensor Programs}

Large language model inference is often presented as a sequence of decode
steps. In agentic systems, those decode steps are embedded in a dynamic control
program. The model may emit a structured tool request, the runtime may validate
the request against a schema, retrieval may fetch relevant memory, and the
result may alter the continuation path. These operations are fine-grained and
dependency-heavy.

The mismatch with conventional GPU execution is not only control flow. It is
also ownership. Tool state may live in CPU memory, vector indexes may live in
GPU memory, parsed arguments may cross host-device boundaries, and scheduling
metadata may be maintained by a runtime that has no cheap way to expose
fine-grained futures to GPU work. As a result, the system pays for orchestration
around the tensor kernels.

\subsection{Agentic Prefill Is Repeated Work}

The usual prefill/decode split treats prefill as the bulk parallel stage and
decode as the latency-sensitive sequential stage. Many current accelerator
efforts therefore specialize decode: minimize per-token latency, optimize
KV-cache reads, and split attention from feed-forward computation so each can be
served by more specialized hardware.

HATS is aimed at a different, under-specialized part of the loop. Tool-using
agents repeatedly create new context that must be validated, retrieved,
filtered, ranked, transformed, and ingested before the model can continue. This
is prefill-like work, even when it occurs in the middle of a conversation. A
tool result may be short, but the system around it may touch resident memory
indexes, schema state, ranking kernels, context buffers, and KV-cache update
paths. HATS names this workload \emph{agentic prefill}.

Agentic prefill has a better PPA shape than general decode control. It exposes
parallelism over retrieved documents, candidate memories, structured fields,
and context blocks. It is often close to model state and local indexes. It can
benefit from coherent memory, cheap queues, and fine-grained signals without
asking the accelerator to own high-level planning or arbitrary external tools.

\subsection{The PPA Question}

The performance, power, and area question is the core discipline for HATS. CPUs
are excellent at irregular scalar control. A new accelerator feature is not
justified simply because it can execute tool-related instructions. It is
justified only if it saves enough latency or energy to compensate for its area
and design complexity.

The strongest PPA case for HATS comes from avoiding unnecessary boundaries:

\begin{itemize}[leftmargin=*]
  \item avoiding CPU wakeups for small, frequent tool operations;
  \item avoiding host-device synchronization when prefill or re-prefill state is
        already GPU-resident;
  \item avoiding copies between model state, retrieval indexes, and tool
        buffers;
  \item keeping structured validation near constrained decoding and context
        ingestion;
  \item batching or fusing tool operations that would otherwise be tiny kernels;
  \item letting a low-area scheduler manage dependency metadata close to the
        queues it serves.
\end{itemize}

The weak case is equally important. HATS should not accelerate arbitrary
browser control, network APIs, shell commands, or general-purpose planning
logic. These operations have variable latency, security-heavy semantics, and
often poor locality with model state. The CPU and operating system should remain
responsible for them.

\section{Design Principles}

\subsection{Specialize Prefill, Not General Decode}

HATS should be understood as a prefill-side accelerator substrate. It is not a
replacement for specialized decode engines, and it should not compete with
hardware whose only job is to minimize one-token latency. Instead, HATS
specializes the preparation of model continuation: retrieval, argument
validation, result filtering, context assembly, placement, and dependency
signals that make the next prefill or decode step cheaper to start.

This distinction keeps the architecture disciplined. The decode path remains
thin and latency-oriented. The HATS path is wider and locality-oriented: it
handles the repeated construction of context around tool boundaries, where
there is enough parallelism and memory traffic to justify heterogeneous
mechanisms.

\subsection{Bounded Tools First}

HATS limits its initial scope to local, bounded tool primitives:

\begin{enumerate}[leftmargin=*]
  \item \textbf{Retrieval}: vector search, key-value memory lookup, embedding
        cache access, and local index traversal.
  \item \textbf{Structured parsing}: JSON parsing, schema validation, argument
        extraction, tokenization, and constrained decoding predicates.
  \item \textbf{Context construction}: packing retrieved or tool-produced data
        into continuation buffers, updating prefill metadata, and preparing
        KV-cache construction paths.
  \item \textbf{Graph scheduling}: futures, dependency tracking, fan-out/fan-in,
        retry, timeout, and cancellation.
  \item \textbf{Data movement}: DMA and placement between model state, tool
        buffers, indexes, CPU memory, storage, and network-facing buffers.
  \item \textbf{Sandboxed pure compute}: deterministic kernels for filtering,
        ranking, scoring, formatting, or lightweight transformations.
\end{enumerate}

This scope is narrow enough to make a hardware conversation meaningful, but
broad enough to cover common hot paths in agentic systems.

\subsection{CPU as Policy, HATS as Local Mechanism}

HATS does not eliminate CPU orchestration. Instead, it changes the boundary.
The CPU remains responsible for high-level planning, OS policy, authentication,
network security, external tool mediation, and arbitrary software execution.
HATS handles frequent local mechanisms near model execution: dispatch, signals,
local retrieval, validation, context construction, and bounded compute.

\subsection{Compiler/Runtime Before Silicon}

HATS should begin as a compiler and runtime project. The first version can lower
agentic tool graphs onto existing CPUs, GPUs, and heterogeneous runtimes. That
implementation provides the measurement basis for hardware proposals. If vector
retrieval dominates, near-memory search engines may be justified. If parsing
dominates, schema units may be justified. If synchronization dominates,
hardware futures and queues may be justified.

\section{HSA Lineage}

HSA introduced a useful model for heterogeneous cooperation. Rather than treat
the GPU as a device controlled exclusively by host driver calls, HSA emphasizes
shared virtual memory, user-mode queues, and low-cost signals between heterogeneous
agents. HATS borrows these ideas while changing the target workload.

\begin{table}[h]
\centering
\begin{tabular}{lll}
\toprule
\textbf{HSA Concept} & \textbf{HATS Adaptation} & \textbf{Agentic Role} \\
\midrule
Shared virtual memory & Coherent tool/model address space & Avoid copies and remapping \\
User-mode queues & Tool dispatch queues & Cheap small-tool launch \\
Signals & Futures and cancellation & Fine-grained dependencies \\
Heterogeneous agents & Tool engines and schedulers & Route work by locality \\
Memory regions & Capability-scoped buffers & Bound tool access \\
\bottomrule
\end{tabular}
\caption{HSA concepts adapted for HATS.}
\end{table}

HATS should avoid reviving older portable instruction set ideas directly.
Instead, it can use modern compiler infrastructure. A plausible stack is an
MLIR dialect for agentic tool graphs, lowering passes to GPU kernels, CPU
tasks, SPIR-V-like compute targets, or sandboxed bytecode, and runtime packets
that look like HSA queues and signals.

\section{Programming Model}

At the programming level, an agentic tool operation can be represented as a
future-producing call with explicit capabilities and schemas:

\begin{lstlisting}
future<Result> tool_call(
    capability cap,
    tool_id tool,
    schema_id input_schema,
    void* input,
    signal cancel
);
\end{lstlisting}

This abstraction is intentionally restrictive. The capability describes the
memory regions and tool classes available to the call. The schema identifies
the expected input structure. The cancellation signal lets speculative or stale
agent branches stop consuming resources. The future becomes a dependency object
that can wake a model continuation, another tool, or a CPU policy handler.

At a lower level, the runtime may represent a tool request as a packet:

\begin{lstlisting}
struct hats_packet {
    uint32_t opcode;          // retrieve, validate, rank, dma, sandbox, ...
    uint32_t capability_id;   // bounds memory and tool access
    uint64_t input_ptr;
    uint64_t output_ptr;
    uint64_t schema_or_kernel;
    uint64_t completion_signal;
    uint64_t cancellation_signal;
    uint32_t flags;
    uint32_t priority;
};
\end{lstlisting}

The packet form matters because it is the bridge between compiler IR and
heterogeneous execution. Different backends can consume the same logical packet:
a CPU task runner, a GPU kernel, a retrieval engine, a parser unit, or a
simulator for proposed hardware.

\section{Compiler and Runtime Architecture}

\subsection{Agent Tool Graph IR}

HATS begins with an Agent Tool Graph IR. This IR captures tool-producing model
steps, dependencies, structured data contracts, and cancellation edges. Unlike a
pure tensor IR, it must represent dynamic control and future-like values:

\begin{lstlisting}
%args = hats.decode_until_tool %model_state
%ok   = hats.validate %args schema(@SearchArgs)
%hits = hats.retrieve %index, %args.query if %ok
%rank = hats.sandbox @rerank(%hits, %args)
%ctx  = hats.assemble_context %rank, %args
%kv   = hats.prefill_context %model_state, %ctx
hats.resume_decode %model_state, %kv
\end{lstlisting}

The IR should be high-level enough to preserve agentic structure and low-level
enough to expose placement decisions. The compiler can decide whether
validation should run as CPU code, GPU code, a parser primitive, or fused logic
near constrained decoding, and whether context construction should be batched
with the next prefill stage.

\subsection{Lowering}

The lowering pipeline can proceed in stages:

\begin{enumerate}[leftmargin=*]
  \item \textbf{Graph normalization}: identify tool boundaries, schemas,
        dependency edges, cancellation scopes, and retry regions.
  \item \textbf{Placement}: choose CPU, GPU, retrieval, parser, DMA, or sandbox
        targets based on data locality and profiling.
  \item \textbf{Packetization}: lower tool operations to HATS packets and
        dependency signals.
  \item \textbf{Fusion and batching}: group tiny operations when their
        dependencies and memory regions permit.
  \item \textbf{Backend emission}: generate CPU tasks, GPU kernels, sandbox
        bytecode, or simulator instructions.
\end{enumerate}

\subsection{Runtime}

The runtime maintains queues, capability tables, signal objects, and memory
regions. Its job is not to make arbitrary distributed systems calls. Its job is
to keep local tool work near the data it consumes and the continuation it
unblocks.

The runtime should expose profiling hooks for:

\begin{itemize}[leftmargin=*]
  \item host wakeups avoided or incurred;
  \item queue latency and signal latency;
  \item bytes moved across CPU-GPU boundaries;
  \item tool operation size distribution;
  \item cancellation latency;
  \item occupancy of retrieval, parser, and sandbox engines;
  \item continuation latency from tool completion to usable prefill or KV state.
\end{itemize}

\section{Candidate Hardware Primitives}

HATS does not require all of the following primitives. They are hypotheses to
evaluate through the compiler/runtime prototype.

\subsection{Queue and Signal Fabric}

A low-area queue and signal fabric may provide the highest return. Agentic tool
work often has many small dependent tasks. If every task requires a host driver
round trip, the system loses even when the computation itself is cheap. Hardware
queues, futures, and cancellation signals can reduce latency and CPU wakeups.

\subsection{Retrieval Blocks}

Retrieval is a strong candidate for near-data acceleration. Embedding indexes,
attention caches, and memory stores may already live in HBM or unified memory.
Vector distance computation, top-k selection, and coarse filtering map naturally
to parallel hardware. The PPA case is strongest when the index is resident and
the result feeds directly into agentic prefill or GPU-side continuation.

\subsection{Structured Parsing and Schema Units}

Agentic systems frequently produce structured arguments. JSON validation,
schema checks, token-level constraints, and format repair can become hot paths
around decoding and context ingestion. Specialized parsers do not need to be large. Even modest
support for finite-state validation, bracket matching, type checks, and
field-level constraints may reduce CPU intervention and failed tool calls.

\subsection{Sandboxed Microkernel Engine}

Some tools are small deterministic functions: filter these hits, rank these
items, transform this structure, compute a score, or normalize this result. A
WASM-like or eBPF-like sandbox target could support these functions with
explicit capabilities and bounded execution. The goal is not general software
compatibility. The goal is safe, predictable local compute near the data.

\subsection{Data Movement Engine}

Many agent workflows are limited by placement rather than arithmetic. A HATS
data movement engine would understand tool buffers, model continuations,
prefill buffers, retrieval indexes, and CPU-visible memory regions. It would use capabilities to
move only authorized regions and signals to wake dependent work.

\section{Security and Isolation}

Tool use is a security boundary. HATS must make this boundary explicit.

Capabilities should specify:

\begin{itemize}[leftmargin=*]
  \item which memory regions a tool can read or write;
  \item which queues it can submit to;
  \item which signals it can wait on or complete;
  \item which schemas or kernels it can invoke;
  \item whether it can trigger CPU-mediated external tools.
\end{itemize}

External tools should remain behind CPU and OS policy. HATS can prepare,
validate, rank, or format requests, but network APIs, browser automation, shell
commands, and filesystem side effects require explicit mediation. This boundary
keeps the first HATS design tractable and prevents the architecture from
becoming a confused general-purpose OS bypass.

\section{Evaluation Plan}

The first HATS prototype should answer whether the substrate is worth building.
It should compare a baseline CPU-orchestrated agent runtime against a HATS-style
runtime on existing hardware.

\subsection{Benchmarks}

Useful benchmarks include:

\begin{enumerate}[leftmargin=*]
  \item \textbf{Retrieve-prefill-resume}: model emits retrieval requests, local
        vector search runs, retrieved context is packed, and decode resumes
        from the resulting prefill or KV state.
  \item \textbf{Schema-heavy tool use}: model emits structured arguments that
        must be validated and repaired before tool dispatch or context
        ingestion.
  \item \textbf{Speculative fan-out}: multiple retrieval or ranking branches
        launch and stale branches are cancelled before re-prefill.
  \item \textbf{Small-kernel pipeline}: many small deterministic transforms run
        between decode and re-prefill steps.
  \item \textbf{Memory locality stress}: tool buffers, context buffers, and
        indexes vary across CPU memory, GPU memory, unified memory, and
        storage-backed regions.
\end{enumerate}

\subsection{Metrics}

The evaluation should report:

\begin{itemize}[leftmargin=*]
  \item end-to-end task latency;
  \item tokens per joule for tool-using decode and re-prefill;
  \item time from tool result to usable continuation context or KV state;
  \item host wakeups per completed tool call;
  \item CPU time per agent step;
  \item bytes copied across host-device boundaries;
  \item queue and signal overhead;
  \item cancellation latency;
  \item failed or malformed tool-call rate;
  \item sensitivity to tool granularity and batch size.
\end{itemize}

\subsection{PPA Modeling}

After profiling the runtime, HATS can model candidate hardware primitives. The
model should be conservative: area estimates for queue fabric, parser units,
retrieval blocks, and sandbox engines should be compared against measured CPU
energy, synchronization overhead, and data movement cost. The useful question is
not whether HATS can beat a CPU on scalar control. The useful question is
whether small local mechanisms can avoid enough boundary cost to justify their
area.

\section{Expected Contributions}

HATS aims to contribute:

\begin{enumerate}[leftmargin=*]
  \item a scoped definition of hardware-friendly agentic tool primitives;
  \item an Agent Tool Graph IR for tool-producing model workflows;
  \item a HATS packet and runtime model inspired by HSA queues and signals;
  \item a placement and lowering strategy for CPU, GPU, retrieval, parser, DMA,
        and sandbox targets;
  \item a benchmark suite for agentic prefill and outer-loop overhead;
  \item a PPA-driven method for deciding which hardware primitives are worth
        adding to future GPUs or heterogeneous packages.
\end{enumerate}

\section{Limitations}

HATS is deliberately narrow. It does not claim that arbitrary tools belong on
the GPU. It does not claim that accelerator hardware should own high-level
agent planning. It does not solve distributed systems latency or make external
APIs deterministic. It also does not assume that every agent workflow benefits
from local tool acceleration.

The HATS hypothesis is strongest when:

\begin{itemize}[leftmargin=*]
  \item model state, indexes, and tool buffers are already local or resident;
  \item tool calls are frequent and fine-grained;
  \item continuation latency matters;
  \item agentic re-prefill events are frequent;
  \item structured validation sits in the prefill/decode loop;
  \item cancellation and dependency scheduling are hot paths.
\end{itemize}

The hypothesis is weakest when tool latency is dominated by remote services,
human interaction, OS side effects, or broad scalar planning.

\section{Conclusion}

Agentic AI exposes a gap between dense model execution and the heterogeneous
tool workflows wrapped around it. Current systems often run the model on a GPU
and the agentic outer loop on a CPU, paying synchronization and data movement
costs at every boundary. While many accelerators specialize the decode path,
HATS specializes the agentic prefill path: the repeated construction,
validation, retrieval, filtering, and ingestion of context before model
continuation. HATS proposes a narrow alternative: keep the CPU in charge of
policy and external side effects, but provide a coherent GPU-adjacent substrate
for bounded local tool primitives.

The project should begin as a compiler and runtime system. That path makes the
hardware story empirical. By representing tool workflows as graphs, lowering
them to queues and signals, and profiling execution on existing systems, HATS
can identify which primitives deserve hardware support. The eventual GPU is not
the starting assumption. It is the consequence of measured agentic prefill
bottlenecks.

\begin{thebibliography}{9}

\bibitem{hsa}
HSA Foundation.
\newblock \emph{HSA Platform System Architecture Specification}.
\newblock Version 1.2, 2017.

\bibitem{mlir}
Chris Lattner, Mehdi Amini, Uday Bondhugula, Albert Cohen, Andy Davis,
Jacques Pienaar, River Riddle, Tatiana Shpeisman, Nicolas Vasilache, and
Oleksandr Zinenko.
\newblock MLIR: Scaling Compiler Infrastructure for Domain Specific
Computation.
\newblock In \emph{CGO}, 2021.

\bibitem{spirv}
Khronos Group.
\newblock \emph{SPIR-V Specification}.

\bibitem{wasm}
World Wide Web Consortium.
\newblock \emph{WebAssembly Core Specification}.

\bibitem{ebpf}
Linux Kernel Documentation.
\newblock \emph{eBPF Documentation}.

\end{thebibliography}

\end{document}

